\section{Evaluation Design}
\label{sec:evaluation}
Organize the evaluation around questions that test the contribution.
For each question, specify the subjects, comparison, metric, and analysis.

\subsection{Questions and experimental controls}
\textbf{RQ1: Task quality.} Does the method improve the target metric under
matched data, tuning budgets, and stopping conditions?
\textbf{RQ2: System cost.} What latency, memory, and compute costs accompany
that change?
\textbf{RQ3: Mechanism.} Which component explains the observed behavior,
and under what conditions does it stop helping?

Report data collection, filtering, splits, and leakage prevention. Record
versions, seeds, hardware, and precision settings. Specify whether
uncertainty reflects independent runs, resampled inputs, or variation
across subjects; these sources answer different questions.

\begin{table}[htbp]
  \caption{A results-table skeleton. Dashes indicate unfilled cells, not
  measured zeroes. Replace them with results and uncertainty estimates;
  do not describe improvements until measurements are available.}
  \label{tab:results}
  \centering\small
  \begin{tabularx}{\linewidth}{@{}Xrrr@{}}
    \toprule
    \textbf{Configuration} & \textbf{Quality $\uparrow$} &
      \textbf{Latency (ms) $\downarrow$} & \textbf{Memory (GB) $\downarrow$} \\
    \midrule
    Reference implementation & \placeholder & \placeholder & \placeholder \\
    Strong baseline          & \placeholder & \placeholder & \placeholder \\
    \method                  & \placeholder & \placeholder & \placeholder \\
    \method{} without constraint & \placeholder & \placeholder & \placeholder \\
    \bottomrule
  \end{tabularx}
\end{table}

\subsection{Results and interpretation}
Use \autoref{tab:results} to report effectiveness and cost. Define the quality
metric. For latency, specify batch size, warm-up, synchronization, and the
measured percentile. Distinguish peak allocated memory from reserved or
resident memory.

Answer each research question and bound its scope. Pair aggregates with
failure analysis. Ablations should change one decision at a time; when
factors change together, describe the confound.
